\section*{Results}\label{section: result}
\subsection*{Charting spatial ligand-target interactions using Renoir}

In order to build a comprehensive computation framework to investigate spatial cell-cell interaction we developed Renoir. As shown in Figure 1, Renoir is an end-to-end computational framework for exploring the spatial map of ligand-target activities either by integrating spatial transcriptomics and scRNA-seq datasets from the same tissue or by employing single-cell resolution spatial transcriptomics data. Renoir first curates a set of ligand-target pairs using NATMI's ConnectomeDB \citep{Hou2020} and NicheNet \citep{Browaeys2020}. For a low-resolution spatial transcriptomic dataset (e.g., 10X Visium), using a well-annotated scRNA-seq data, Renoir quantifies a neighborhood activity score for each curated ligand-target pair for each spot in the spatial transcriptomic data. For a single-cell resolution spatial transcriptomic data, Renoir quantifies the neighborhood activity score for each curated ligand-target pair for each cell with a spatial context. The quantification of the neighborhood activity score for a specific ligand-target pair for a particular spot is performed using the proportions and cell type-specific mRNA abundances present in each spot/cell within the defined neighborhood of the spot/cell (see Methods for details). The activity between a ligand-target pair for the given spot and another spot in the neighborhood is scored based on the cell type–specific mRNA abundances of the two genes weighted by the cell type proportions as well as the mutual information between the two genes across the cell types present within the two spots. Thus, Renoir generates a novel neighborhood activity score matrix (of dimensions number of ligand-target pairs $\times$ number of spots). Using the ligand-target neighborhood activity score matrix, Renoir performs several downstream analyses. First of all, Renoir infers spatial communication domains through unsupervised clustering of spots based on ligand-target neighborhood activity scores and determines communication domain-specific ligand-target activities. Moreover, Renoir can infer cell type-specific ligand-target activities within a communication domain. Renoir further provides an algorithm for ranking the ligands within a communication domain based on their cumulative activity across target genes expressed by the major cell types in the domain. Renoir also provides a mapping of pathway-level activity across spots and communication domains. Finally, Renoir encompasses a visualization module for the visualisation of ligand-target activity across spots as well as a Sankey plot that provides an overview of user-specified ligand-target activities across cell types and communication domains.

\subsection*{Renoir outperforms state-of-the-art methods in inferring ligand-target interactions}
\textcolor{blue}{Filled out citations and missing values. }We first evaluated Renoir's performance in inferring the spatial activity of ligand-target pairs using simulated datasets generated using a semi-synthetic enrichment-based simulation workflow (Figure 2a, see Methods for details). Specifically, we started with paired single-cell and spatial datasets based on which we simulated new spatial datasets where activity of a ligand-target pair was simulated in reference spatial locations (positive reference spots) through the enrichment of cell types associated with the ligand and target gene. Moreover, the simulated datasets also contained spatial locations which did not harbor any activity between the ligand and target (negative reference spots). To comprehensively cover the range of interactions across tissue types, we used three reference matched scRNA-seq and ST datasets from human intestine \citep{FAWKNERCORBETT2021810}, triple negative breast cancer (TNBC) \citep{Wu2021} and dorsolateral prefrontal cortex (DLPFC) \citep{Maynard2021} for the simulation. The simulation workflow as described in Figure 2a was applied to generate 43, 15 and 26 simulated datasets corresponding to the human intestine, TNBC and DLPFC datasets respectively, where, each simulated dataset enriched a specific ligand-target interaction. Since Renoir quantifies ligand-target activity at a specific spatial context and is the only method (to the best of our knowledge) to do so contrary to the other cell-cell interaction methods, which infer ligand-receptor interactions by employing the spatial data, as competitor methods, we selected state-of-the-art ligand-receptor interaction inference methods COMMOT \citep{Cang2023}, STLearn \citep{Pham2023} and SpatialDM \citep{Li2023}. These three methods are capable of quantifying ligand-receptor interactions at the resolution of spatial spots and we adopted them for quantifying ligand-target activity from our simulated datasets. The performance of each method was evaluated based on a spatial activity accuracy metric that computes the Pearson correlation of the spot-level ligand-target activity score measured by a method with the true simulated activity of the ligand-target pair in the positive and negative reference spatial locations (see Method for details). In addition, we also compared the methods in terms of precision and recall (see Methods) for the ligand-target pairs to evaluate their ability to distinguish the spots harboring ligand-target interactions from the spots devoid of such interactions.\par 

In each simulated dataset across the tissue types, Renoir was able to capture the ligand-target interactions more precisely. To illustrate, we consider the ligand-target pair \textit{A2M-MYC}, simulated from the human intestine dataset (Fig. 2b), where the positive (blue) and negative (red) reference spots included enrichment and abatement of said ligand-target pair interaction respectively as is also evident from the spatial expression plot of the two genes. Amongst the positive spots, we observed Renoir's ligand-target activity scores to truly reflect the enriched spots with high scores as compared to COMMOT and stLearn. While SpatialDM had high scores for positive spots, it also had scores all across the tissue and failed to differentiate the positive and negative spots appropriately. Renoir was also able to correctly quantify the absence of ligand-target interactions in the negative reference spots. In comparison, all three competitor methods showed false-positive activity of the gene pair in multiple negative spots (red bounded regions labeled 1, 2 and 3) with near obsolete \textit{A2M} and \textit{MYC} expression. Renoir is capable of computing the activity of a ligand-target by considering the expression of a suitable receptor at the receiving cell. To showcase this, we considered the pair \textit{IFNG-JCHAIN}, simulated using the TNBC dataset, where \textit{IFNG} was associated with Natural Killer T Cells (NKT Cells) and CD8$^{+}$ T Cells; and \textit{JCHAIN} was associated with B Cells and Dendritic Cells (DCs). However, B cells do not contain a receptor for \textit{IFNG} and thus the \textit{IFNG} secreted from NKT Cells or CD8$^{+}$ T Cells can not activate \textit{JCHAIN} in B cells. Therefore, the positive reference spots harbored interactions with DCs as the target cell type whereas the negative spots contained the enrichment of B cells as the target cell type (see methods section for more details). In this scenario, by accounting for the presence of suitable receptors while calculating ligand-target activity scores, Renoir correctly scored the negative spots with negligible scores; whereas other methods failed to consider such scenarios and therefore incorrectly gave high scores to the negative spots as depicted by the red bounded regions (Supplementary Figure 1a). Similarly, for the DLPFC dataset (e.g., for the pair \textit{NPY:CARTPT}) also, Renoir's spatial scores for the ligand-target activity were more precise as compared to the other methods (Supplementary Figure 1b). SpatialDM and COMMOT had less distinction between the positive and non-positive spots and stLearn failed to capture the activity pattern due to the high range of \textit{NPY} expression and relatively low range of \textit{CARTPT} expression. Quantitatively, Renoir outperformed other CCC inference methods in inferring spatial ligand-target activity by achieving 45.8, 98.5, and 12.4\% \textcolor{blue}{[calculated as so: average across (renoir median)/(method median) - 1]} improvement in the median spatial activity accuracy across the human intestine, TNBC and DLPFC datasets (Fig 2c) indicating Renoir's superiority in capturing the spatial variation of the ligand-target activity. Furthermore, in terms of average precision and recall over the ligand-target pairs, Renoir outperformed SpatialDM, COMMOT, and stLearn across the tissue types across different threshold values used for calculating precision and recall (Supplementary Figure 1c). For varying threshold values, Renoir was able to consistently maintain high precision and recall values across the tissue types. In comparison, the precision and recall values for the other methods varied across datasets, with a method having high precision-recall for one dataset but low precision-recall for another. \par 

Finally, we evaluated the accuracy of the spatial communications domains inferred based on the ligand-target interaction scores inferred by a CCC inference method. For this, we utilized 10X Visium datasets from 12 dorsolateral prefrontal cortex (DLPFC) samples \citep{Maynard2021}, in which the authors manually annotated six cortical layers and white matter for each sample based on cytoarchitecture and selected marker genes. Considering the manual annotations as ground truth, we evaluated the utility of the ligand-target scores inferred by Renoir in identifying the layers and compared its performance against that of COMMOT, SpatialDM and stLearn. Each method was evaluated based on adjusted Rand index (ARI) and normalized mutual information (NMI) scores calculated by comparing the communication domains inferred by the method against the ground truth annotations over all 12 DLPFC samples (see Methods for details). Renoir achieved the highest ARI and NMI scores considering all the samples and outperformed all the other methods (Fig. 2d) indicating that it is able to capture spatial variation and architecture better than the other CCC inference methods. Fig. 2e shows how the communication domains inferred by Renoir captured the underlying spatial structure for the sample 151675 better than the other methods for which the inferred communication domains were much noisier.


% We also evaluated the possible spatial distinctiveness of the communication domains inferred through the scores generated by each of the method. We can evaluate the methods with respect to their ability to uphold the underlying spatial structure (if available) within their ligand-target interaction scores. To do so, we consider the DLPFC dataset. Maynard et. al have manually annotated each of the spots as white matter or DLPFC layer. Considering this as the ground truth, we calculated the NMI/ARI scores amongst the inferred communication domains for each method and the ground truth annotations across all available DLPFC samples (see methods for details). Renoir holds the most spatial information in comparison to the other methods (Fig \textcolor{red}{5d}) with Renoir showing the best average NMI/ARI scores. This is can be seen from the spatial plots of communication domains vs ground truth where Renoir has communication domains that are closest to the underlying spatial structure.\\

% In order to compare the methods over their ability to make a distinction between spots enriched with cell type interactions and spots that are not, we calculate the precision and recall of the method in correctly deducing a spot as a positive reference spot or negative reference spot (see method section for details). We compute the average precision and recall for each method; across the simulated datasets for each ligand-target pair considered in the Human Intestine, TNBC and DLPFC datasets. We observe that for varying score thresholds, Renoir is able to maintain high precision and recall values in comparison to the other methods (Supp Fig 12c). This implies that the distinction between scores generated by Renoir for positive and negative reference spots are more clear cut in comparison to other approaches. Furthermore, Renoir is able to consistently maintain high precision and recall scores. As for the other methods, the Precision and Recall values tend to vary across datasets, with a method having high precision/recall for one dataset and low precision/recall for another.
% In terms of precision-recall as well, Renoir outperformed SpatialDM, COMMOT, and stLearn across the tissue types across different threshold values used for calculating precision-recall (Supplementary Figure c). \par 

% The simulation follows enrichment of cell types associated with \textit{A2M} (Epithelial cells) and \textit{MYC} (Pericytes). This is showcased in Figure 5b,    \\ 
% Two simulated samples are generated such that each sample has a common ligand cell type with different target cell types as shown in Supp Fig 12a. In both the samples, we perform simulation with the assumption that the B Cells do not have a suitable receptor to interact with the NKT Cells or CD8+ T Cells. 

% Ligand-target interactions can occur only if a suitable receptor is expressed at the receiving cell. In order to consider this possibility, we consider the ligand-target pair \textit{IFNG-JCHAIN}, simulated across the TNBC dataset. Here \textit{IFNG} is associated with Natural Killer T Cells (NKT Cells) and CD8+ T Cells; and \textit{JCHAIN} is associated with B Cells and Dendritic Cells (DCs). Two simulated samples are generated such that each sample has a common ligand cell type with different target cell types as shown in Supp Fig 12a. In both the samples, we perform simulation with the assumption that the B Cells do not have a suitable receptor to interact with the NKT Cells or CD8+ T Cells. Therefore, the positive reference spots are cell type interactions with DCs as the target cell type and the negative spots are cell type interactions with B Cells as the target cell type (see methods section for more details). In this scenario, Renoir accounts for the presence of suitable receptors while calculating ligand-target activity scores and therefore correctly scores the negative spots with negligible scores; whereas other methods fail to consider such scenarios and therefore incorrectly score the negative spots well as depicted by the red labelled bounded regions. To further showcase the precision of Renoir, we consider the ligand-target pair \textit{NPY:CARTPT}, simulated across the DLPFC dataset (Supp Fig 12b). Here the ligand and target cell types enriched are Somatostatin inhibitory neurons (IN-SST) and NRGN-expressing neurons (Neu-NRGN-I) respectively; where the positive and negative spots are simulated similar to as in the Human Intestine dataset. We can observe that Renoir shows a more clear scoring gradient than other methods across the positive and negative spots. SpatialDM and COMMOT fail to provide a clear scoring gradient amongst the positive and non-positive spots. On the other hand, stLearn fails to capture activity patterns due to the high range of \textit{NPY} expression and relatively low range of \textit{CARTPT} expression. Overall, Renoir outperforms other CCC inference methods in inferring simulated cell type interactions across the x, y and z simulations performed across the Human Intestine, TNBC and DLPFC datasets, as is depicted by the box plot in Fig 5c.\\

% We performed a comparative analysis of Renoir's performance in inferring cell type interactions against prominent cell-cell communication (CCC) inference methods over simulated datasets as showcased in the simulation workflow (Figure 5a). More specifically, cell type specific ligand-target interactions are simulated over existing spatial transcriptomic datasets over which scores are generated using Renoir and three robust CCC inference methods namely COMMOT, STLearn and SpatialDM (permutation and z-score) all of which represent the diverse methodologies involved in CCC inference. The ligand-target activity scores generated are then compared against the simulated cell type interaction scores to find the method with the best ability to deduce the simulated cell type interactions (see methods section). We simulate three datasets, Human Intestine, sample X of the TNBC Cancer and Dorsolateral Prefrontal Cortex (DLPFC). We apply the simulation workflow as described in Figure 5a to generate x, y and z simulated datasets for the Human Intestine, TNBC and DLPFC datasets respectively. Here, each of the simulated dataset focuses on a given cell type specific ligand-target interaction as explained in the methods section. In each of the datasets, we observed Renoir to capture cell type specific interactions more precisely. 


% We also evaluated the possible spatial distinctiveness of the communication domains inferred through the scores generated by each of the method. We can evaluate the methods with respect to their ability to uphold the underlying spatial structure (if available) within their ligand-target interaction scores. To do so, we consider the DLPFC dataset. Maynard et. al have manually annotated each of the spots as white matter or DLPFC layer. Considering this as the ground truth, we calculated the NMI/ARI scores amongst the inferred communication domains for each method and the ground truth annotations across all available DLPFC samples (see methods for details). Renoir holds the most spatial information in comparison to the other methods (Fig \textcolor{red}{5d}) with Renoir showing the best average NMI/ARI scores. This is can be seen from the spatial plots of communication domains vs ground truth where Renoir has communication domains that are closest to the underlying spatial structure.\\
% In order to compare the methods over their ability to make a distinction between spots enriched with cell type interactions and spots that are not, we calculate the precision and recall of the method in correctly deducing a spot as a positive reference spot or negative reference spot (see method section for details). We compute the average precision and recall for each method; across the simulated datasets for each ligand-target pair considered in the Human Intestine, TNBC and DLPFC datasets. We observe that for varying score thresholds, Renoir is able to maintain high precision and recall values in comparison to the other methods (Supp Fig 12c). This implies that the distinction between scores generated by Renoir for positive and negative reference spots are more clear cut in comparison to other approaches. Furthermore, Renoir is able to consistently maintain high precision and recall scores. As for the other methods, the Precision and Recall values tend to vary across datasets, with a method having high precision/recall for one dataset and low precision/recall for another.

% Intracellular functional gene interactions are computed and weighted using mutual information between two genes. Two intra- cellular networks are connected via cross-talk edges (i.e., known ligand-receptor interactions). Ligand-receptor pairs with higher cell type–specific gene expression but lower correlated expression with- in the same cell type (thus more likely to be involved in cross-talk instead of self-talk) are assigned higher cross-talk weights. Nodes in the integrated network are weighted by a combination of their cell type–specific gene expression and closeness to the ligand/receptor genes in the network.

% In order to generate a ligand-target pair vs spots "neighbourhood score" matrix, Renoir uses (i) scRNA-seq data for a given sample and its corresponding cell type annotations, (ii) corresponding spatial transcriptomic data, and (iii) a curated list of ligand-target pairs as shown in Figure [reference]. As can be seen in Figure [reference], each cell of the neighbourhood score matrix scores the ligand-target activity that can be found in the neighbourhood of a certain spot (for more information, see Methods). The cell type distribution across the spatial topology and the cell type-specific mRNA abundance across each spot within the spatialtranscriptomics data can be determined with the help of cell2location \citep{Kleshchevnikov2022}. NicheNet \citep{Browaeys2020} is utilized to predict over a set of secreted ligand-receptor pairs available in NATMI's ConnectomeDB in order to generate a list of curated ligand-target pairs (see Methods). Calculations of prior cell type specific ligand-target activity are performed using the scRNA-seq data available (see Methods). An information-theory-based scoring metric is applied over each and its corresponding neighbourhood in order to determine the value of each cell within the neighbourhood score matrix which is illustrated in Figure [reference]. The neighbourhood scores that have been generated can be used for a variety of downstream tasks, as are illustrated in Figure [reference]. These tasks include (i) visualisation of ligand-target activity across spots; (ii) inference of the cell type, ligand, and target specific communication domains; (iii) determination of communication domain specific ligand-target activity; (iv) mapping geneset/pathway activity across spots and communication domains; (v) a ranking of ligand specific activity across communication domains; and (v) a Sankey plot that provides an overview of ligand-target activity across each cell type and communication domain.

\subsection*{Renoir identifies distinct communication domains in Mouse brain}

To demonstrate Renoir's ability to uncover region-specific ligand-target interactions, we first applied Renoir to 10X Visium data from two Mouse brain sections (samples S1 and S2) and its corresponding single nuclei data \citep{Kleshchevnikov2022} (see data availability). Renoir identified $9409$ (S1) and $9405$ (S2) ligand-target pairs that were expressed in this dataset and quantified their spatial neighborhood activity scores for both samples. Based on the ligand-target neighborhood activity scores, Renoir clustered the ST spots across the two tissue sections and identified $19$ spatial communication domains which had distinct latent space embeddings (Figure 3a) and were similarly distributed across the sections (Figure 3b). We further quantified the distribution of the spatial communication domains across different anatomical brain regions and found that the domains across both samples primarily corresponded to the same anatomical regions and were highly concordant across the two samples in terms of their distribution in anatomical brain regions (Figure 3c).\par

The communication domains were found to harbor diverse cell types that were spatially localized and abundance pattern concordant across the replicates (Supplementary Figure 2a-b). The domains further harbored differential activity of specific ligand-target pairs (Figure 2d, Supplementary Figure 3b-c). For instance, domain 1 specific to white matter exhibited \textit{Il33} and \textit{Csf1} activity which are known to be required for microglia development and maintenance in white matter \citep{10.3389/fimmu.2019.02199, 10.3389/fimmu.2018.02596} and domain 3 specific to the hypothalamus exhibited \textit{Dlk1} activity known to be expressed in hypothalamic arginine vasopressin and oxytocin neurons \citep{10.1371/journal.pone.0036134}. Domains 2 and \textcolor{blue}{10} located in the thalamus region were enriched in \textit{Tcf7l2} and \textit{Lef1} activity which are known to be essential for neurogenesis in the developing mouse neocortex \citep{Chodelkova2018, Nagalski2013}. We also noticed the existence of ligand or target-specific activity pertaining to distinct anatomical regions, wherein the ligand or target was known to be abundant within the brain regions corresponding to the domain. For example, domain \textcolor{blue}{9} exhibited high activity of \textit{Rspo1}, which is known to be highly expressed within the cortex \citep{10.1210/en.2013-1550, Miller2019}, domain \textcolor{blue}{7} exhibited high activity of hippocampal marker \textit{Crlf1}, \textcolor{blue}{domain 6 exhibited high \textit{Tgfa} activity, where TGF-$\alpha$ is synthesised across regions which included the tuber olfactorium, a part of the ventral striatum \citep{Wilcox1901},} and domain \textcolor{blue}{13} exhibited high activity of \textit{Ccnd1} and \textit{Nr2f2}, which are known to be highly expressed within the Amygdala region \citep{KOELLER2008230, 10.1242/dev.075564}. Supplementary Table 1 lists the domain-specific ligand-target interactions with literature support for the associated brain region. 

%Furthermore, we also observed section-specific interactions that were specific to S1 or S2. S2 specifically showed \textit{Tgfa} activity within domain 7 which was located mainly in the striatum region. This is in line with previous findings of . \par 


% We first generated communication domains over the integrated neighborhood scores across the two sections (see methods section for more details). We observed that the domains generated amongst the two sections were similarly distributed across the sections as seen in figure [reference]. To validate our observations, we cross tabulated spots labelled with respect to the domains they belong to, to the available region annotations of each spot, as described in the methods section. We acquire a frequency table as shown in figure [reference], where the domains across both samples primarily corresponded to the same regions and the distribution of domains across the two sections can be seen to be similar. 

% The communication domains were found to harbor spatially localized cell types and ligand-target activity as can be seen in figure [reference], where the cosine similarity amongst spots were found with respect to the cell type distribution and ligand-target activity. As a consequence, each domain consisted of differentially expressed ligand-target activity, as seen in figure [reference].  

Kleshchevnikov et al. \citep{Kleshchevnikov2022} identified and molecularly characterized 10 regional astrocyte subtypes using this spatial and matched single-cell data. However, interactions of these regional astrocyte subtypes with other spatially co-localized cell types were not explored. We applied Renoir to characterize the interaction of these regional astrocyte subtypes with other cell types. Figure 3d and Supplementary Figure 3d illustrate the interactions where ligands expressed by regional astrocyte subtypes activate other target genes in other cell types across samples S1 and S2 respectively. Figure 3e and Supplementary Figure 3e demonstrate the interactions where ligands expressed by other cell types activate specific marker genes in astrocyte subtypes. Ligands expressed by the telencephalic astrocytes in Cortex, and Amygdala were found to activate genes in excitatory neurons in Layer 6b (\textit{Vegfa:Ctgf}) and Amygdala (\textit{Ptn:Nr2f2}) respectively. On the other hand, \textit{Agt} expressed by diencephalic astrocytes (\textcolor{blue}{Thalamus}) were found to activate genes in inhibitory neurons (\textcolor{blue}{\textit{Psap:Pde4dip}}). White matter astrocytes interacted with oligodendrocytes via \textit{Tnc} which activated gene \textit{Adamts4}. Similarly, different genes in telencephalic astrocytes in Amygdala, Cortex, and Hippocampus were found to be activated by ligands expressed by excitatory neurons in Layer \textit{4} (\textit{Rspo1:Vegfa}), Amygdala (\textcolor{blue}{\textit{Slit1:Nr2f2}}), and Hippocampus DG (\textcolor{blue}{\textit{Crlf1:Ccnd3}, \textit{Slit1:Tcf4}}) respectively. Excitatory neurons in the thalamus were found to activate genes in diencephalic astrocytes \textcolor{blue}{(Thalamus Medial (\textit{Ntng1:Gpld1}, \textit{Nxph1:Tcf7l2}))}.

%Lateral Thalamus astrocytes that were abundant in the ventrolateral thalamic nuclei, expressed \textit{Serpine2} which activated different genes in different excitatory neurons in the thalamus.

% We further examined region specific astrocyte subtype ligand-target activity across the two samples (see methods section). Figure [reference] involves activity where the ligands are astrocyte subtype markers and figure [reference] involves activity where the targets are astrocyte subtype markers. We observe localization of multiple ligand-target interactions specific to associated regions of the astrocyte subtypes. For instance, we can see activity ($Tgfa:Gfap$, $Apoe:Mfge8$, $Slit1:Tcf4$) driven by subtype markers $Gfap$, $Apoe$ and $Tcf4$ of subtypes White Matter, Amygdala and Hippocampus respectively and cell type markers $Tgfa$, $Mfge8$ and $Slit1$ of cell types Oligodendrocyte Progenitor Cells (OPC 2), Amygdala astrocyte and Excitatory Hippocampus (Ext HPC) respectively. \textcolor{orange}{I've found the following literature which denotes some form of relation amongst the mentioned pairs, but I'm not sure if they're relevant here: \\
% https://www.sciencedirect.com/science/article/abs/pii/016538069500026A?via\%3Dihub,\\ https://www.sciencedirect.com/science/article/abs/pii/S000689930703017X?via\%3Dihub,\\ 
% https://www.ncbi.nlm.nih.gov/pmc/articles/PMC5018780/,\\ https://www.sciencedirect.com/science/article/pii/S2589004222013748}\\

For each communication domain, based on the agreement of neighborhood activity and prior regulatory potential scores, we ranked the prominent ligands that are expressed by the major cell types in the domain (see Methods for details). As shown in Figure 3f and Supplementary Figure 3f, the top-ranked ligands for some of the major domains with literature support for sections S1 and S2 respectively. This further highlighted the region-specific enrichment of certain ligands. For instance, \textcolor{blue}{\textit{Bdnf} known to be associated with numerous functions within the hippocampus \citep{Heldt2007} showcases high activity within domain 16 associated with the Hippocampus. \textit{Nrg1}, known to have high impact on cortical functions \citep{AGARWAL20141130} exhibited high activity within domain 5.} Supplementary Table 2 lists such major ligands with their corresponding literature support.

% We also ranked prominent ligands within each communication domain with respect to a neighborhood score and prior ligand target  (see Methods section). On doing so, we were able to observe ligand specific activity known to be abundant within the domain it is associated with. For instance, we observed $Slit1$ known to be extensively expressed within the hippocampus [reference], prominently active in domain 17 associated with the hippocampus.   

Using Renoir, we spatially mapped major pathway activities based on the activity scores of the ligand-target pairs associated with the pathways (see Methods). We found the pathway activity to be concentrated in brain regions that are typically associated with the said pathway (Figure 3g, Supplementary Figure 3g and Supplementary Figure 3). For example, the cortex and hippocampus were highly enriched in MAPK signalling and NEUROTROPHIN signalling pathways, respectively. This is consistent with prior works which demonstrated high concentrations of MAPK components in the cortex and hippocampus of the adult mouse brain \citep{https://doi.org/10.1002/cne.21186}, and NEUROTROPHIN signaling's role in the regulation of neurogenesis \citep{10.3389/fnins.2016.00026}. 

\subsection*{Renoir uncovers intratumor spatial communication sub-domains within triple negative breast cancer}

To demonstrate Renoir's ability to characterize complex heterogeneous tissues in terms of ligand-target interactions, we analyzed single-cell and Visium datasets (Figure 4a) from a Triple Negative Breast Cancer (TNBC) patient \citep{Wu2021}. From the set of curated ligand-target pairs, Renoir identified $18654$ ligand-target pairs that were expressed in this dataset and quantified their spatial neighborhood activity scores. Based on the neighborhood activity scores of ligand-target pairs, Renoir clustered the spatial spots into six communication domains which were also distinct in low-dimensional embeddings (Figure 4b). While each of the spatial communication domains was primarily associated with morphologically distinct pathological regions, some communication domains were distributed across multiple pathological regions (Figure 4c). The domains were further differentiated based on the ligand-target activities and the composition and abundance of the cell types they harbored (Figure 4d-e, Supplementary Figure 7a-b). Since two (clusters 4 and 5) of the six domains had low overall UMI count (Supplementary Figure 6c), we focused on the first four communication domains in our further analyses. In accordance with their respective cell type compositions, interesting differential ligand-target activities were observed across the six domains. For instance, domain 0 which was primarily associated with the histological region annotated as ductal carcinoma in situ (DCIS) and secondarily contained spots pathologically annotated as stroma and lymphocytes, harbored a high abundance of cancer cycling and cancer basal SC cell types (Figure 4c, Supplementary Figure 7a). Moreover, it also harbored cancer-associated fibroblasts (CAFs) and different T cell populations (Supplementary Figure 7a). Domain 0 demonstrated differential activity of \textit{MYC} (e.g., \textit{LAMB1-MYC}, \textit{SERPINE2-MYC}), \textit{MDK} (\textit{MDK-CCNA2}) and \textit{SAA1} (\textit{SAA1-ESRRA}) which were all marker genes for the cancer cell types in this dataset (Figure 4d-e). While \textit{MYC} and \textit{SAA1} are known to be overexpressed in basal-subtype of breast cancer \citep{doi:10.1073/pnas.191367098,doi:10.1073/pnas.1732912100,OT26566}, \textit{MDK} is a known cancer progression marker \citep{Filippou2020}. Furthermore, domain 0 also harbored high activity of \textit{VEGF} (\textit{VEGFA}, \textit{VEGFB}) (known to be highly up-regulated in breast cancer \citep{RIBATTI2016453,Harmey1998,Lewis2000}), and other stromal ligands (\textit{COL18A1}, \textit{COL4A1}) (Figure 4e). Domain 1 was associated with the histological regions lymphocytes and invasive cancer and majorly harbored cancer-associated myofibroblasts (myCAF), macrophages, and T cells (cycling and \textit{CD8}$^{+}$) along with the two cancer cell types (Figure 4c, Supplementary Figure 7a). Consequently, interactions associated with tumor progression and invasion (e.g., \textit{CXCL10-MMP2}, \textit{CCL3-ICAM1} \citep{doi:10.1073/pnas.1408556111}, \textit{FN1-MMP10} \citep{MCGOWAN20081566}), as well as suppression (e.g., \textit{CXCL9-IRF1} \citep{10.1158/1078-0432.CCR-19-1868}) among these cell types were enriched in domain 1 (Figure 4e). Domain 2 was associated with the histological region normal, stroma and lymphocytes and harbored four major cell types including inflammatory-like CAFs (iCAF-like) and myoepithelial cells (Figure 4c, Supplementary Figure 7a). Interestingly, this domain harbored high activity of the ligand \textit{CCL28} which is known to promote breast cancer and metastasis \citep{Yang2017}, target \textit{HES1} which plays a crucial role in malignant cell proliferation \citep{Li2018} and \textit{AREG} which is expressed by myoepithelial cells and loss of which can inhibit tumor proliferation, growth and invasiveness \citep{Mao2018} (Figure 4e). Domains 3 and 4, despite low UMI count, showcased distinct ligand-target activities. For example, domain 3 displayed distinct interleukin activity (e.g., \textit{IL6}, \textit{IL33}, \textit{IL1RN}) which are known to play important roles in metastasis \citep{ARA20101223},\citep{10.3389/fimmu.2014.00141},\citep{Shiiba2015} and \textit{PENK} was found to be differentially active in domain 4 (Figure 4e). \par 

% Several macrophage and myCAF secreted ligands had Different tumor growth inhibitors secreted by CAFs and macrophages were found to be active within domain 1 (e.g., TGFB, CCL3, CXCL9, etc.). Other than communications between suppressants [reference] (Eg: CCL3-IRF1), communications between tumor suppressing \citep{Liu-Chittenden2017} and tumor promoting \citep{Zhang2020} \citep{Chen2022} genes were found (Eg: RARRES2-SERPINE1) in domain 1. 

We further investigated the most prominent cell type-specific ligand-target interactions (Figure 4f, Supplementary Figure 8a) within each spatial communication domain (see Methods for details). By emphasizing domains with the highest ligand-target activities and UMI count (0, 1, 2, and 3), we uncovered cell type-specific interactions such as \textit{RARRES2-AEBP1}, \textit{CXCL12-GSN} (ligand and target are markers of iCAFs, domain 2). \textcolor{blue}{We also noticed cross-communication between myCAFs and \textit{CD4}$^{+}$ T cells (\textit{CXCL12-FASLG}) in domain 1 [it was FASLG instead of ITK]}. Tumor metastasis and progression-promoting interactions between macrophages and cancer cycling cells (\textit{CXCL16-CCND1}) \citep{OT3690, cancers13143568} were observed in domain 3. Additionally in domain 3, we observed interaction between \textit{TGFB1} and \textit{FOXP3} which are markers of NK cells and \textit{CD4}$^{+}$ T cells respectively indicating \textit{TGFB1}-mediated suppression of immune response \citep{10.1371/journal.pone.0076379, doi.org/10.1111/imm.12941}. \par  
% \citep{Zarzynska2014}.\par 
% We also noticed high CXCL9 activity within clusters 1 and 3 (CXCL9-MMP2, CXCL9-MMP9) which are abundant in immune cell types, which correlates with the fact that CXCL9 plays an important role in modulating the infiltration of immune cells into the Breast Cancer microenvironment \citep{10.3389/fonc.2021.710286}. 

For each communication domain, we further ranked the major ligands (expressed by the major cell types in the domain) according to their activities (see Methods for details). A key EMT marker, \textit{TGFB3} \citep{jcm5070063} was one of the top-ranked ligands in both domains 0 and 1 (which harbored a high abundance of cancer cells) and its top targets included \textit{FN1} (in myCAFs and macrophages) and \textit{TNC} (in myoepithelial cells and myCAFs) which are associated with breast cancer metastasis \citep{Wang2018}, \citep{Oskarsson2011} (Figure 4g). In domain 1 which contained high abundance of macrophages, another top ligand was \textit{CCL2}, which is known to promote metastasis via retention of metastasis-associated macrophages (MAMs) \citep{10.1084/jem.20141836} and it was found to target \textit{ICAM1} (a molecular target for TNBC) \citep{doi:10.1073/pnas.1408556111} in myoepithelial cells and macrophages, other chemokines \textit{CCL3} and \textit{CCL4} as well as lymphocyte inhibitory molecule \textit{LTB} in \textit{CD4}$^{+}$ T cells (Figure 4g). Within the top ligand lists of each communication domain, we also found ligands that were specifically active in a single communication domain (Figure 4h). For example, in domain 1, \textit{CD8}$^{+}$ T cells secreted the ligand \textit{CCL5} (known to promote TNBC progression through the regulation of immunosuppressive cells \citep{Zhang2013}) which further activated \textit{MMP9} and \textit{CCL2} in macrophages. \par

Subsequently, utilising cell type abundance values of spatial locations, we sub-clustered each communication domain (domains 0, 1, 2, and 3) into subdomains that were unique in terms of cell type abundance and entailed distinct ligand-target activities (Supplementary Figures 9-12). Amongst the \textcolor{blue}{six} subdomains in domain 0 (Supplementary Figure 9), in subdomain \textcolor{blue}{4} \textit{TGFB1} and \textit{FBN1} activated several EMT-associated genes (\textit{FBN1}, \textit{MMP2}, \textit{COL5A1}, \textit{PDGFRB}, etc.) \citep{10.1172/JCI45014} and cancer progression and regulation markers (\textit{DUSP1}, \textit{JUNB}, and \textit{MYC}) \citep{Tuglu2018, Lukey2016} respectively. In contrast, subdomain \textcolor{blue}{1} which harbored different immune cells (\textit{CD4}$^{+}$ and \textit{CD8}$^{+}$ T cells, DCs, NKT cells) as well as iCAFs, showcased extensive chemokine activity (\textit{CCL5}) which targeted different genes in different immune cells and iCAFs (e.g., \textcolor{blue}{\textit{EGR1}}). 


Domain 1 was segregated into \textcolor{blue}{4} distinct subdomains (Supplementary Figure 10), of which subdomain \textcolor{blue}{2} abundant with immune cells\textcolor{blue}{[should be Macrophages here now, since sub domain 3 with RARRES2 activity has high macrophage and myCAF abundance but not much immune cells]} showcased high \textit{RARRES2}  activity (secreted by myCAFs) and subsequently a lower abundance of cancer cells compared to other subdomains. This is coherent with previous findings which show that \textit{RARRES2} serves as an immune-dependent tumor suppressor by acting as a chemo-attractant to recruit anti-cancer immune cells \citep{wittamer2003specific,wittamer2005neutrophil,zabel2005chemokine,vermi2005role}. The subdomain \textcolor{blue}{2}, which harbored a high abundance of macrophages, myoepithelial cells, myCAFs, and cycling myeloid cells displayed \textit{FBN1} and \textit{IL1B}-mediated activation of \textit{MMP1}, which is known to promote brain metastasis in TNBC \citep{Liu2012}. Additionally, subdomain \textcolor{blue}{1} abundant with cancer cell types showcased high activity of \textit{TGFA}, a ligand associated with tumor invasion \citep{doi.org/10.1002/ijc.28295} and \textit{DKK1-HES1}, which are known to promote tumor metastasis, proliferation and invasion \citep{zhu2021expression, Li2018}. High \textit{DLK1} (a prominent marker for iCAFs \citep{10.1158/2159-8290.CD-19-1384}) activity was observed in subdomain \textcolor{blue}{3} of domain 2 (Supplementary Figure 11) as it was abundant with iCAFs. Interestingly, the only subdomain (subdomain \textcolor{blue}{0}) of domain 2 that harbored cancer basal cell type showed \textit{NTN4}- (expressed by myoepithelial and mature luminal cells) and \textit{SEMA3C}- (expressed by mature luminal cells) mediated activity of \textit{FOXA1} indicating potential roles of these co-localized cell types in EMT-progression \citep{anzai2017foxa1} (Supplementary Figure 11). Domain 3 was segregated into five subdomains (Supplementary Figure 12) of which subdomain \textcolor{blue}{4} abundant with immune cells showed high \textit{CCL19} \textcolor{blue}{\sout{and \textit{CCL21}}} activity which are known to attract lymphocytes \citep{vicari2000antitumor} and subdomain \textcolor{blue}{0} abundant with cancer cells, myCAFs, and myoepithelial cells showcased activity of \textit{SERPING1} (metastasis-predictor extracellular matrix gene \citep{doi:10.1073/pnas.0710331105}), \textit{TGFA}, \textit{COL4A1}, and \textit{SPP1} respectively.
% we observed activity between NTN4 and FOXA1 in subdomain 3 of domain 2 located in the Normal + stroma pathologically annotated region, which includes scarce cancer cells and abundant epithelial cells. NTN4 expression is known to induce N-cadherin downregulation \citep{Xu2017} and FOXA1 is known to promote E-cadherin expression , which also aligns with the high luminal epithelial concentrations within the subdomain that is responsible for expressing E-cadherin \citep{daniel1995expression}.
% What is interesting to note is the co localization of TGFB1 and FBN1 activity and additionally the interaction between TGFB1 and FBN1 which aligns with the importance of FBN1 in EMT \citep{LIEN2019375}. 
 
% \citep{ueno2000significance} \citep{luboshits1999elevated} \citep{silzle2003tumor} \citep{owen2005expression}. 

Using Renoir, we further spatially mapped the major pathway activities based on the activity scores of the ligand-target pairs associated with the pathways (see Methods). While some pathways showed activity across multiple communication domains, some other pathways were more active in specific communication domains (Supplementary Figure 8b). For example, the WNT signaling pathway was found to be active in subdomain 1 of domain 0, and subdomain 2 of domain 1 which harbored a high abundance of cancer cell types (Figure 4i). IL-6/JAK/STAT3 signaling was highly active in subdomain 1 of domain 1 (Figure 4i) that harbored a high abundance of myCAFs, macrophages and myoepithelial cells along with cancer cells but lacked immune cells other than cycling T-cells indicating a highly immunosuppressive tumour microenvironment \citep{Johnson2018}. TGF beta signaling was active in domains 0 and 1, and MAPK signaling had the highest activity in domain 2 (Figure 4i). Taken together, Renoir was able to provide comprehensive insights into the spatial heterogeneity of ligand-target activity and its implication in tumor biology.


\subsection*{Renoir identifies hepatocyte-macrophage interactions in developing fetal liver}

Cell-cell communication shapes tissue functioning and organogenesis during development. In order to evaluate the versatility of Renoir in uncovering cell-cell interactions during development, we applied Renoir to data from fetal liver development. We generated a 10X Visium spatial transcriptomic dataset (see Methods for details) from a fetal liver sample for which single-cell RNA sequencing data was available \citep{SHARMA2020377} (Figure 3a). Renoir identified $17630$ ligand-target pairs active in the tissue and quantified their spatial neighborhood activity scores. Based on the ligand-target neighborhood activity scores, Renoir identified five spatial communication domains (Figure 3b), which differed based on ligand-target activity and distribution of UMI count across the tissue section (Figure 3b-c and Supplementary Figure 4a). Domain 0 demonstrated the majority of ligand-target activity, followed by domain 1. On the other hand, domains 2 and 4 showed sparse activity (UMI counts were low) (Figure 3c) while domain 3 was specific to the target \textit{CDKN1A} (Supplementary Figure 4a). This is in line with the UMI count distribution across spots, where the UMI count was most abundant in domains 0 and 3, followed by domains 1, 2 and 4. Cell types were found to be evenly distributed across domains 0, 1 and 3 (Supplementary Figure 4b). Activities pertinent to fetal liver were found to be widespread (Figure 3d). For example, \textit{COL18A1} associated with type XVIII collagen (a prominent component of the Extracellular Matrix (ECM) in the liver \citep{10.1242/dmm.011577}) was found to be highly expressed across the tissue and activate multiple targets including hematopoietic transcription factors \textit{BCL2L1} and \textit{BCL6} \citep{10.1371/journal.pone.0007641}. Higher \textit{CDKN1A} activity was observed, in line with the proliferative nature of hepatocytes in developing fetal liver tissue.\par
Further analyses involved acquiring previously described ligand rankings and prominent cell type-specific ligand-target interactions for each communication domain. Specifically, we observed prominent activity of the ligands \textit{MDK}, an important growth factor for fetal liver hematopoietic stem cells and multipotent progenitor expansion \citep{Gao2022} and \textit{CXCL12} which helps in retaining hematopoietic stem cells in the fetal liver \citep{Lewis2021} (Figure 3e). We further observed the occurrence of ligand-target pairs \textit{CXCL12-CXCR4} and \textit{HGF-COL1A2} as prominent interactions (Supplementary Figure 5a). While \textit{CXCL12} is known to interact with \textit{CXCR4} to control the regeneration of liver amongst other organs \citep{10.3389/fimmu.2020.02109}, \textit{CXCR4} is known to be expressed at the highest density by the lymphoid cells \citep{https://doi.org/10.1002/SICI}. Moreover, the deletion of one of \textit{CXCL12} or \textit{CXCR4} is known to cause defective B-cell lymphopoiesis \citep{10.3389/fimmu.2020.02109}. Importantly, we also observed \textit{MARCO}$^+$ tissue-resident macrophage population (Kupffer cells) interacting with \textit{APOB}$^+$ hepatocytes. \textit{MARCO} (MAcrophage Receptor with COllagenous structure) is predominantly expressed in non-inflammatory Kupffer cells \citep{MacParland2018} and fetal-like reprogramming of liver macrophages plays an important role in immunosuppressive phenotype in liver cancer \citep{SHARMA2020377}. As shown in Figure 3e and Supplementary Figure 5a, we observed that \textit{MARCO}$^+$ macrophages interact with \textit{PLG}$^+$/\textit{APOB}$^+$ subset of hepatocytes. The \textit{PLG} (plasminogen) has been known to be important for hepatic regeneration \citep{doi:10.1073/pnas.96.26.15143,00001721-200809000-00006} and \textit{PLG} deficiency can lead to impaired remodelling in mouse models after acute injury. Next, we subdivided hepatocytes into two groups: \textit{ALB}$^+$\textit{APOA}$^+$\textit{PLG}$^+$ and \textit{ALB}$^+$\textit{APOA}$^+$\textit{PLG}$^-$ (Figures 4f-g). As shown in Figures 4g-h, these two subsets are transcriptionally and spatially distinct from each other indicating heterogeneity of the fetal liver hepatocyte population. We observed that the majority of the \textit{ALB}$^+$\textit{APOA}$^+$ hepatocyte population does not express \textit{PLG} while there is a specific cluster with a high expression of \textit{PLG} (Figure 3f). In the Visium data, the spots expressing \textit{ALB}$^+$\textit{APOA}$^+$\textit{PLG}$^+$ hepatocytes also contained \textit{FOLR2}$^+$ macrophages (Figure 3h). \textit{FOLR2}$^+$ macrophages are embryonic macrophages present in the fetal liver microenvironment \citep{SHARMA2020377}. Spatial correlation analyses further showed higher similarity of \textit{FOLR2}$^+$ macrophages and \textit{ALB}$^+$\textit{APOA}$^+$\textit{PLG}$^+$ hepatocytes as compared to \textit{PLG}$^-$ hepatocytes in terms of cell type distribution as well as neighborhood activity score (Figure 3i).\par
Major pathway activities were also spatially mapped onto the Fetal Liver sample (Supplementary Figure 5c). Although domain specific activities were undetected, relevant activity notably, KRAS signaling (plays an important role in Fetal Liver hematopoiesis \citep{10.1002/stem.2355}) and canonical JAK STAT signaling (recognised to be integral in initiating a signalling cascade required for embryonic development, haemopoietic development and differentiation \citep{doi:10.3109/08977194.2012.660936}) were distinctly observed. Taken together, these results provide important insights into the macrophage-hepatocyte niche in the developing human liver.

\subsection*{Renoir identifies onco-fetal interactions in hepatocellular carcinoma}
To evaluate Renoir's ability in handling single-cell resolution ST datasets, we applied it on a Nanostring CosMx dataset generated from a hepatocellular carcinoma (HCC) patient. The dataset consisted of x cells distributed across y cell types and y genes. 

Two recent studies discovered the presence of onco-fetal tumor microenvironment (TME) in hepatocellular carcinoma, where the \textit{FOLR2}$^{+}$ TAMs, \textit{PLVAP}$^{+}$ ECs and \textit{POSTN}$^{+}$ CAFs were identified as onco-fetal cell types that shared transcriptional programs with their fetal liver counterparts. We wanted to evaluate Renoir's ability in identifying the onco-fetal interactions from the Nanostring CosMx dataset. To perform an unbiased analysis, we used Renoir to quantify the activity of the ligands expressed by the onco-fetal cell types with the other major cell types present in the neighborhoods of the onco-fetal cells using Renoir's ligand-ranking formulation (Figure 6a). Renoir was able to characterize the activity of several known core onco-fetal ligands (e.g., \textit{VEGFA}, \textit{TGFB3}) as well as ligands shared by two onco-fetal cell types (e.g., \textit{BMP7}, \textit{CXCL12}, \textit{TGFB1}, \textit{IFNG}) as depicted in Supplementary Figure b. Specifically, Renoir identified the activity of known onco-fetal interactions \textit{CXCL12}-\textit{CXCR4} and \textit{VEGFA}-\textit{KDR} (Figure 6c). CXCL12 was expressed by \textit{PLVAP}$^{+}$ ECs and \textit{FOLR2}$^{+}$ TAMs and it activated CXCR4 in \textit{POSTN}$^{+}$ CAFs as well as macrophages and Tregs. Renoir correctly identified the co-localization of these cells and quantified the activity of the pair in the neighborhood. Similarly, VEGFA was expressed by all three onco-fetal cell types and it interacted with \textit{KDR} in co-localized ECs and \textit{MUC6}$^{+}$ bipotent hepatocytes (Figure 6b). Among the cell types co-localized in the onco-fetal niche, Renoir identified the \textit{MUC6}$^{+}$ bipotent hepatocytes to have the highest number of interactions with the onco-fetal cell types (more specifically with \textit{POSTN}$^{+}$ CAFs) (Figure 6c). We delved into interactions specific to oncofetal cells and \textit{MUC6}$^{+}$ bipotent cells by quantifying the cell type-specific activity of the relevant ligand-target pairs in the onco-fetal neighborhoods (Supplementary Figure c). We noticed the activation of MYC and TP53 in \textit{MUC6}$^{+}$ bipotent cells by multiple onco-fetal ligands including \textit{BMP2, TGFB3, CXCL12}. We utilized gprofiler \citep{gProfiler} to infer the  pathways in the \textit{MUC6}$^{+}$ bipotent cells targeted by the oncofetal ligands and identified various interleukin signaling pathways (e.g., IL4, IL13, IL10), cytokine signaling and prostaglandin signaling (Figure 6d) to be enriched. Since prostaglandin signaling is known to promote HCC progression \citep{prostaglandin_pathway}, we further looked into the activity of \textit{IL6-PTGS2}. High scores for \textit{IL6-PTGS2} was found in the tissue niche where \textit{MUC6}$^{+}$ bipotent cells were co-localized with the oncofetal cells indicating the potential role of onco-fetal cells in the activation of PTGS2 in the \textit{MUC6}$^{+}$ bipotent cells. 

% To do so, we selected ligand-target pairs with high average scores in neighborhoods harboring \textit{POSTN}$^{+}$ CAFs and \textit{MUC6}$^{+}$ bipotent cells (Supp Fig ).
%% domain details
% We generated Renoirs scores and clustered them into 6 domains \textcolor{red}{[not sure if domains is the right word here since we don't have spatial domains]} (Supp Fig ). We noticed that domains 2 and 5 had low UMI counts (Supp Fig) and therefore decided to utilize domains 0,1,3 and 4 for further analysis. We noticed that domains x and y harboured \textit{TP53} and \textit{MYC} specific interactions. On looking into the cell type distribution plot (Supp Fig) we noticed high proportions of A and B cell types across domains x and y. From this, we further wanted to evaluate Renoir's ability in capturing the onco-fetal interactions in the TME from CosMX data. 

% Additionally, DLK1 based activity were found to be widespread in regions abundant in iCAFs (sub domain 2 of domain 2) and is known to be a marker for the same \citep{10.1158/2159-8290.CD-19-1384}}.\par

% We applied LTSpatial over the TNBC dataset (scRNA-seq data, corresponding cell type annotations and spatial transcriptomic data) to gain the neighborhood scores over the set of curated ligand-target pairs available. The neighborhood scores were utilized to generate spotwise-interaction specific clusters as seen in figure [reference]. We noticed that clusters corresponded to six distinct domains (figure [reference]) where each domain was differentiated on the basis of ligand-target activity, cell types and UMI counts across the tissue section as described in figures [reference]. With respect to UMI count, the domains are broadly distinguished into two groups of high UMI count (which include domains 0, 1 and 2) and low UMI count (which include domains 3, 4 and 5). The cell type distribution across the six domains can be seen in figure [reference].  For instance, domain 0 which was abundant specifically in Cancer cycling and Cancer Basal cell types amongst others, showcased ligand-target activity which included high VEGF (VEGFA, VEGFB) activity which are known to be highly up-regulated in breast cancer [reference]. It also involved SERPINE activity which are known to support the creation of tumor-promoting conditions [reference]. Domain 1 is abundant in cancer-associated myofibroblasts (myCAF), macrophages, T cells (CD4+ and CD8+) amongst other cell types. High number of tumor suppressants were found to be active within domain 1 (TGFB, CCL3, CXCL9 etc.). Other than communications between suppressants [reference] (Eg: CCL3-IRF1), communications between tumor suppressing [reference] and tumor promoting [reference] genes were found (Eg: RARRES2-SERPINE1). Domain 2 was found to be abundant in inflammatory-like CAFs (iCAF-like) amongst other cell types. Interestingly, high activity of the ligand CCL28 was found within domain 2 which is known to promote breast cancer and metastasis [reference] amongst others such as HES1 which is crucial in malignant cell proliferation [reference] and AREG, loss of which is known to stunt tumor proliferation, growth and invasiveness [reference].  

 % Further analysis allowed us to peer into cell type specific ligand-target activity across each domain (see Methods for more details), where domain prominent ligand-target interactions were observed such that the ligand and target were markers of specific cell types. Interestingly enough, we noticed cell type specific interactions such as RARRES2-AEBP1, CXCL12-GSN where the ligand and target are markers of iCAFs at domain 2 and HEBP1-CCND1 amongst Cancer Basal cells present in domains 0 and 1 to name a few. We also noticed cross communication between Cancer Basal cells and Cancer cycling cells (MDK-CCND1, MDK-RAN), CD4+ T cells and iCAFs (CXCL12-CXCR4, CXCL12-ITK), CD8+ T cells with Cancer Basal and Cancer cycling cells (CXCL13-CCND1, GZMB-CCND1) to name a few within domains 0 and 1. Next we performed our pathway analysis to map groups of ligand-target interactions, particularly the available gene sets and generated DeNovo clusters as shown in figure [reference]. Domains 0, 1 and 3 were further clustered with respect to the cell type abundance across the domain to generate 5, 3 and 2 ``sub domains" respectively as shown in figures [reference] with their respective cell type distributions as shown in figures [reference]. Amongst multiple gene sets depicted in figures [reference], we observed the localisation of pathways such as MAPK signalling in domains 2 and 1 (sub domain 0 and 2) and WNT signalling in domain 0 (sub domain 0) respectively, where there is a presence of Cancer cell types (iCAFs are predominant in domain 2 and sub domain 0 of domain 1, myCAFs are predominant in sub domain 0 of domain 1 and Cancer cycling cells are predominant in sub domain 0 of domain 0). This aligns with previous observations, wherein MAPK is known to play a crucial role in the survival and development of tumour cells [reference] and WNT is known to be associated with breast cancer [references]. Further more, it is interesting to note the presence of pairs such as [need to mention pairs here]. [one liner on De Novo cluster here]. [Interesting points from sankey plot]



\begin{comment}
To showcase the versatility of LTSpatial across varying datasets, we analysed a [time duration regarding dataset] Fetal Liver dataset as depicted in figure [reference]. We first applied LTSpatial over the provided Fetal Liver scRNA-seq data, its corresponding cell type annotations and spatial transcriptomic data to gain the neighborhood scores over the set of curated ligand-target pairs available (see section [reference] for more details). The neighborhood scores were utilized to generate spotwise-interaction specific clusters as seen in figure [reference]. We noticed that clusters corresponded to distinct domains where each domain was differentiated on the basis of distinct ligand-target activity and distribution of UMI count across the tissue section. Cluster 0 showcased majority of ligand-target activity, followed by cluster 1. Cluster 2 and 3 showed sparse activity ([mention relevant activity here]) and cluster 4 was specific to the target CDKN1A. This is in line with the UMI count distribution across spots, where the UMI count was most abundant in cluster 0 and 4, followed by cluster 1, 2 and 3. Varying clustering resolutions (see section [reference]) results in distributed ligand/target specific niches are found specific to common ligands or targets (figure [reference]). Further analysis generated cell type-region specific ligand-target interactions (see section [reference] for more details), where domain specific ligand-target interactions were observed such that the ligand and target were markers of specific cell types. Specifically, we observed [need to mention interesting activity here]. Additionally, our pathway analysis showcased multiple signaling pathways across the spatial topography. Majority of the pathways were concentrated across cluster 0 where UMI count is highly dense. Nevertheless, different signalling pathways showcased varying spatial patterns, some of which coincided across the spatial topography and some were distinct from each other, insinuating the possibility of overlapping signalling pathway ``grooves". For instance, figure [reference] depicts the gradient of specific signalling pathways namely JAK STAT, WNT SIGNALLING, MAPK and TGF BETA signalling pathways. Here we observe that while there are highly overlapping pathways like the JAK STAT and TGF BETA, pathways such as WNT and MAPK showcase relatively more distinct spatial localization. Further, the sankey diagram showcases [specific information from sankey plot to be added here].
\end{comment}